\chapter{Conclusion}
\label{ch:conclusion}

This thesis began with a simple question: could reinforcement learning make a simulated booster
land? The first cylinder showed that it could, but it also made the question less interesting.
With an oversized engine, constant mass and no aerodynamics, a successful landing said little
about how the controller would behave when the vehicle changed. The project therefore became an
attempt to build the experiment around the policy, not only the policy itself.

That change in focus explains the final simulator. Vehicle hardware, physical parameters,
objectives, curricula, learning settings, faults and telemetry can all be configured through one
workflow. Editable values become immutable sessions before training starts. Observation and
action dimensions follow the selected hardware, incompatible models are rejected, and every
continuation creates a traceable revision. These mechanisms grew out of practical failures: a
configuration panel made shared state ambiguous, modular hardware broke manually maintained
policy dimensions, and resumed runs needed to preserve rather than overwrite their history.

The three tasks then tested the system at increasing levels of difficulty. The revised fixed-hover
policy completed all 250 \SI{60}{s} horizons, reducing terminal speed to \SI{0.073}{m/s} and
tilt to \SI{0.100}{\degree}, although its \SI{4.77}{m} mean planar offset and the absence of a
binary criterion prevent a success-rate claim. Moving-target hover completed 191 of 250
deterministic trials and remained reliable through $d=0.75$, before failing at the exact
curriculum boundary. The final L11 landing policy completed 212 of 250 trials,
including 18 of 50 at $d=1$. Compared with L10, it reduced mean engine restarts from 18.46 to
1.56 and mean first-contact vertical speed from \SI{3.02}{m/s} to \SI{1.33}{m/s}. The policy
learned navigation, braking, engine sequencing and stable four-foot support, although its 36\%
full-difficulty success shows that the final boundary is not solved reliably.

The L1--L11 campaign explains how that result was reached. A policy that appeared to fly away
revealed a missing spawn velocity. A policy that reached the pad and crashed revealed a capped
impact penalty. Frequent contact revealed that one foot was not a landing, and a policy that
learned with easier actuators revealed excessive engine cycling. Each failure made the next
question narrower. Telemetry and component tests were what turned surprising behaviour into a
specific change to physics, reward, termination, actuation or software ownership.

The central lesson is that an RL policy learns the environment it receives, including its
mistakes. More reward terms do not necessarily express intent more accurately, and more physical
detail does not automatically make an experiment trustworthy. The useful combination is an
explicit model, rewards whose contributions can be inspected, clear ownership of state and an
evaluation that remains separate from training.

The final deterministic landing gap is therefore not an inconvenient result to hide; it is the
next question exposed by the system. A matched stochastic-inference test can determine whether
the continuous representation of binary engine-enable actions contributes to that gap. The
hover task similarly turns its residual planar offset into a specific next test: freeze an
acceptance radius and hold time, then change one suspected cause at a time. Later studies can add
vehicle comparisons, wind and faults under the same discipline. The project has moved from
making one object land to building a product in which the reason it landed---or failed to
land---can be examined.
